% ЗГЕНЕРОВАНО scripts/build_format_v2_chapter.py — не редагувати вручну.
% Джерело: books/ashtavakra-gita/poetychnyi/rozdil-13.md

\poemverse{13.1}{Стан неволодіння — це бути в собі; \\ навіть у жебрацькому поясі він рідкісний. \\ Відкинувши віддавання й приймання, \\ я сиджу, як мені завгодно.}
\poemverse{13.2}{Десь втомлюється тіло, \\ десь втомлюється язик, \\ десь розум. \\ Відкинувши це, \\ я стою в меті людини, щасливо.}
\poemverse{13.3}{Подумавши по-справжньому: \\ зробленого не буде нічого — \\ що б не випало зробити, \\ зробивши те, \\ я сиджу, як мені завгодно.}
\poemverse{13.4}{У йога, що перебуває в тілі, \\ є стани, до яких його змушують \\ дія й недіяння. \\ Та через відсутність зʼєднання й розʼєднання \\ я сиджу, як мені завгодно.}
\poemverse{13.5}{Ні користі, ні шкоди немає мені \\ від стояння, ходіння чи лежання. \\ Тому, стоячи, ходячи чи сплячи, \\ я сиджу, як мені завгодно.}
\poemverse{13.6}{У сплячого немає мені втрати, \\ як і немає досягнення \\ в того, хто докладає зусиль. \\ Відкинувши загибель і піднесення, \\ я сиджу, як мені завгодно.}
\poemverse{13.7}{Побачивши не раз у всіх станах \\ те правило, що має вигляд задоволення, \\ відкинувши добре й погане, \\ я сиджу, як мені завгодно.}
